Трасса исполнения (далее -- трасса) представляет собой упорядоченный набор записей о событиях, происходивших в ходе трассировочного запуска приложения \cite{pinplay,whole_execution_traces}. Подобные трассы широко применяются в индустрии для решения различных задач, таких как исследование поведения приложений, анализ производительности платформ, отладка, сбор профилей исполнения для компиляторов \cite{pinplay,autofdo,lisitsyn_disser}. В рамках данной работы рассмотрены трассы для RISC-архитектуры.

Целевые сценарии исполнения, записанные в виде трасс, могут применяться для анализа поведения разрабатываемых платформ. При этом настройки системы памяти новой платформы (далее -- новая конфигурация или новые настройки) могут отличаться от настроек платформы, на которой были записаны трассы. Возникает необходимость адаптации трасс под новую конфигурацию системы памяти.

Целью данной работы является разработка решения, позволяющего изменять трассы исполнения для модификации записанных в них образов памяти приложений в соответствии с ограничениями новой конфигурации системы памяти.

Представленное решение осуществляет перенос областей образа памяти в соответствии с новыми ограничениями посредством поиска и модификации адресов и смещений в трассе. Для генерации необходимых модификаций анализируется граф потока данных трассы. Рассмотрены алгоритмы генерации карты перемещений областей памяти, сбора графа потока данных из трассы исполнения, поиска локаций адресов и смещений в трассе с помощью обходов графа, обработки граничных случаев с помощью симуляционного прохода по модифицированной трассе. Описанный подход позволяет обеспечить консистентность потоков управления и данных в результирующих трассах, а также снизить влияние модификаций на репрезентативность этих трасс. Представленные алгоритмы могут быть использованы для других задач, связанных с модификацией адресов и смещений в трассах исполнения.

Временная и пространственная сложность обработки трассы линейна по числу инструкций и связей по данным между ними. Предложена метрика для предварительной оценки влияния производимых модификаций на репрезентативность трасс. Рассмотрены другие метрики для количественной оценки свойств реализации. Представлены результаты экспериментальных запусков разработанного средства, результаты валидации выходных трасс и значения собираемых метрик на наборе целевых трасс заказчика для одной из рассматриваемых конфигураций: с увеличенным размером страниц и расширенными областями, зарезервированными под данные ядра ОС. Описаны возможности для дальнейшего улучшения разработанного средства и новые направления его применения.

Предложенное решение успешно применяется для анализа поведения разрабатываемой системы на целевых сценариях исполнения и конфигурациях системы памяти.

\newpage
